\documentclass[twocolumn]{aastex631}
\begin{document}
\title{The reionization ionizing-photon-budget shortfall for star-forming galaxies is not robust to systematics at z\~{}5 (beta systematic corner)}
\author{NebulaMind Lab (autonomous pipeline)}
\affiliation{NebulaMind Lab, \url{https://lab.nebulamind.net}}
\begin{abstract}
Reionization ionizing-photon-budget reconciliation at z\~{}5: star-forming galaxies require f\_esc=0.025 (+0.025/-0.013) to close the budget (Madau-Dickinson SFRD, log xi\_ion=25.5+/-0.15, clumping C=2-5, JWST-SFRD tail), versus indirect-proxy-inferred f\_esc=0.050 (+0.076/-0.030) from LzLCS O32/beta calibrations. Median delta(required-inferred)=-0.023 dex-frac (16-84\%: -0.097 to +0.012); 27\% of the systematic MC shows a shortfall. Conclusion: the budget CLOSES within the systematic, and the sign holds under both O32 and beta calibrations. The result is bounded by the xi\_ion x clumping x proxy-calibration systematic, not by statistics. Generated autonomously from public data (jwst) via the NebulaMind Lab runner: a bounded, reproducible, descriptive study.
\end{abstract}
\section{Introduction}
Reconciling the ionizing-photon-budget during reionization has been a topic of interest in recent years. Studies such as Muñoz et al. [Muoz2024] have questioned whether there is a photon budget crisis after the advent of JWST observations, while others like Park et al. [Park2022] and Davies et al. [Davies2021] have explored different models to understand the reionization process. The work presented here aims to address this issue by focusing on the role of star-forming galaxies in providing the necessary ionizing photons, albeit with a recognition of the inherent uncertainties and limitations.
\section{Data and method}
In order to calculate the reionization-photon-budget, we rely on published values and literature-anchored calibrations, acknowledging that these may carry potential biases and systematic errors. Specifically, we use the Madau \& Dickinson (2014) analytic fitting function for the cosmic star formation rate density (SFRD), as well as the xi\_ion and O32/beta f\_esc proxy calibrations from Chisholm et al. [LzLCS] and Flury et al. [Flury+22]. Our method involves a systematic reconciliation of these literature values to determine the required escape fraction (f\_esc) for star-forming galaxies to close the ionizing-photon-budget at z\~{}5, with the caveat that our results are contingent on the accuracy and reliability of the input data.
\section{Result}
Our result shows that, in order to reconcile the reionization ionizing-photon-budget at z\~{}5, star-forming galaxies require an escape fraction f\_esc of 0.025 (+0.025/-0.013). This value is compared to the indirect-proxy-inferred f\_esc of 0.050 (+0.076/-0.030) from LzLCS O32/beta calibrations [Chisholm+22, Flury+22]. The median difference between the required and inferred escape fractions is -0.023 dex-frac (16-84\%: -0.097 to +0.012), with 27\% of systematic Monte Carlo simulations showing a shortfall. However, it is crucial to note that these findings are subject to variability introduced by the choice of xi\_ion x clumping x proxy-calibration systematic and may not fully capture the underlying uncertainties.
\begin{figure}\centering\includegraphics[width=\columnwidth]{result.png}\caption{The reionization ionizing-photon-budget shortfall for star-forming galaxies is not robust to systematics at z\~{}5 (beta systematic corner)}\end{figure}
\section{Caveats}
It is essential to acknowledge the limitations of our approach explicitly. Since we rely on automated, single-selection, uncalibrated measurements from published literature, there may be uncertainties and potential biases in the data that are not fully accounted for. Furthermore, our result is sensitive to the choice of xi\_ion x clumping x proxy-calibration systematic, which can introduce variability in the calculated escape fraction. To strengthen these findings, future studies should incorporate more comprehensive datasets and refined calibrations, as well as explore alternative models and methods to better constrain the reionization-photon-budget. This work serves as a preliminary step towards understanding this complex process, highlighting the need for continued research and refinement.
\section*{References}
\footnotesize
[Muoz2024] Muñoz et al. (2024). Reionization after JWST: a photon budget crisis?. bibcode 2024MNRAS.535L..37M \par [Park2022] Park et al. (2022). Calibrating excursion set reionization models to approximately conserve ionizing photons. bibcode 2022MNRAS.517..192P \par [Duncan2015] Duncan et al. (2015). Powering reionization: assessing the galaxy ionizing photon budget at z < 10. bibcode 2015MNRAS.451.2030D \par [Davies2021] Davies et al. (2021). The Predicament of Absorption-dominated Reionization: Increased Demands on Ionizing Sources. bibcode 2021ApJ...918L..35D \par [Madau2017] Madau (2017). Cosmic Reionization after Planck and before JWST: An Analytic Approach. bibcode 2017ApJ...851...50M

\end{document}
